%% TODO: Reference -- solving Zeno paradox of state transition
%%       Relational Ontology of Time (Leibniz) vs. Possibility of real time

\documentclass[11pt, a4paper]{article}
\usepackage[utf8]{inputenc}
\usepackage[T1]{fontenc}
\usepackage{geometry}
\usepackage{amsmath}
\usepackage{amssymb}
\usepackage{amsthm}
\usepackage{cite}
\usepackage{titlesec}
\usepackage{fancyhdr}
\usepackage{enumitem}

% Geometry settings
\geometry{margin=1in}

% Title formatting
\title{\textbf{An Approach to The Cosmogony of the Conceptual Universe}}
\author{\textit{Dr.-Ing. Frank-René Schäfer, Independent Researcher}}
\date{\today}

% Custom environments for formalisms
\newtheoremstyle{formalism}% name
  {10pt}%      Space above
  {10pt}%      Space below
  {\itshape}%  Body font
  {}%          Indent amount
  {\bfseries}% Theorem head font
  {.}%         Punctuation after theorem head
  {.5em}%      Space after theorem head
  {}%          Theorem head spec (can be left empty, meaning 'normal')

\theoremstyle{formalism}
\newtheorem{formalism}{Definition}

\begin{document}

\maketitle

%% TODO mention that the physical universe is only an instance of a conceptual 
%% universe that exists only the the framework of conceptualization. This text,
%% tries to find the origin of conceptualization itself. we start from void.
%% Problem: The following explains the evolvement of the universe and a cognitive
%% subject. The contradiction exists with our physical model of consciousness,
%% namely, that consciousness (the cognitive subject) pops into existence after
%% the establishment of time (nitty detail billions of years after that). The only
%% solution to offer is that the universe pops into existence of a cognitive observer
%% at the moment it starts to exists. for the cognitive observer time starts there
%% for him no alternative time can exist. 

\begin{abstract} 
TODO
\end{abstract}

\section*{Nomenclature}

%% TODO

\begin{abstract} 

\end{abstract}

\section{The Noumena and the Cognitive Subject}

We begin with a fundamental demarcation: The \textit{Physical Universe}---the
domain of atoms, neurons, and galaxies---is strictly an instance of a
\textit{Conceptual Universe}. Physics is a discipline of models, and models are
conceptual structures. Therefore, an inquiry into the origin of reality cannot
start with physical time (e.g., the Big Bang), for physical time is itself a
concept that requires a cognitive framework to be intelligible.

We begin with the axiom of \textit{Finite Receptivity}~\cite{Kant1781}. There
is the {\it Cognitive Subject\/} and the {\it Noumena\/}---the latter being the
``Outside'' of the Subject, the unknowable, un-modeled, un-conceptualizable
Absolute Reality. As Kant rigorously demonstrated, the thing-in-itself, the
Noumena, remains forever beyond the grasp of the cognitive subject.

This text proposes a mathematical framework for describing the metaphysical
process underlying a {\it Cognitive Subject\/}'s interaction with absolute
reality. We aim to identify the mathematical mechanics at the heart of our
perception of the universe---mechanics which evade physical observation but
reveal themselves through logical necessity.

TODO:\ I am not sure if this is needed anymore:

In the text, the expression ``faculty of X'' shall serve as a shorthand for the
``ability to experience the existence of X''. We start from a minimum set of
undeniable experiences. The Subject is defined by its exposure to the following
primordial faculties that precede any specific content:

\begin{enumerate}
    \item \textbf{Equivalence:} The faculty to experience a shared 
    characteristic between two distinct entities. 

    \item \textbf{Differentiation:} The faculty to experience a unique 
    characteristic ($\delta$) that distinguishes an entity from another. 

    \item \textbf{Instantiation:} The faculty to experience \textit{Novelty}---the 
    appearance of an entity that has not been experienced before.

    \item \textbf{Immutability:} The faculty to experience an 
    entity to remain in existence.

    \item \textbf{Necessity (Implication):} The faculty to experience the 
    generation of an entity $B$ purely because $A$ exists. 

\end{enumerate}

TODO:\ until here

\section{The Primordial Event and the Origin of Fixation}

We must begin from a conceptual void, where only the Noumena exists: an
un-conceptualizable Absolute Truth---a ground reality---an existence strictly
independent of the conceptualization of a Cognitive Subject. At this origin
point, no fixation, no writing, nothing that could constitute time, is assumed
to exist; there is only the initial fundamental atomic unit of conceptual
nature: an {\it Event\/}---caused by the only existing candidate of action,
namely the Noumena. 

\begin{formalism}[Event]

    An \textbf{Event} $e$ is the atomic instance of a `concept', an
    information, i.e.\ something that can be experienced/perceived, that
    `informs'. It has no temporal expansion. 

    An event has two properties, namely the truth of its own existence $\top$
    and a {\it discernible\/} $\delta$ by means of which it possesses identity.

    \begin{equation}
        e := \langle \top, \delta \rangle
    \end{equation}

\end{formalism}

For the original event, the discernible $\delta_\dot$ has only Identity based
on its being different from the Noumena. It establishes the faculty of {\it
Differentiation}. The only true information possible at this point, must be
about the only existing entity, i.e.\ the Noumena. It is the Noumena's
information of ``I exist''. Due to the inability to interact with any other
entity in the nascent `world of conceptions', $e_\dot$ is unable to develop any
type of dynamic. It is immediately the first converged \textbf{Fixed Point}. It
constitutes the first instance of immutability, stability, and thus, the first
moment of \textbf{Fixation}. Still in a void of anything that may establish any
concept of time, the fixated information $\delta_\dot$ of the Noumena's ``I
exist'' becomes the anchor for all truth to come.

Information only exists if it informs. If an observer possesses information
(knows) about something, this `something' does not inform and is not
information. By necessity, the existence of information $\delta_\dot$ implies
the existence of a lack of knowledge: an entity that can be informed, namely
a {\it Cognitive Subject}.

Once having informed, i.e.\ the event has been fixated, the information has
become knowledge and ceeds to exists---it does not inform anymore. This is the
first grain of dynamic. It establishes the event $e_{:}$ carrying the
information $\delta_{:}$ of the Noumena's ``I instantiate/inform''. Fixation of
the dynamic event into the immutable establishes faculty of a difference
between the {\it Dynamic\/} of {\it Becoming\/} and the {\it Static\/} of {\it
Have-Become\/} in terms of `informing' a cognitive Subject. 

Once the Event has informed the Subject, its function as active \textit{Informing}
is exhausted. It undergoes a \textbf{Phase Transition} from Validity ($\top$)
to Stability ($\bot$). It ceases to be a Signal and becomes a \textbf{Trace}.

This Trace constitutes the \textbf{Tangible Substance} of the Subject's
reality. The Subject does not `remember' the event; the Subject `observes' the
Trace. Thus, the Immutable History is not a mental archive, but the
\textbf{Constituted World} itself. The information sedimented from an event
with a truth modality of $\top$ into an entity with a truth modality of $\bot$
is later defined as the \textbf{Atomic Fact}. It is the matter out of which the
tangible, perceivable universe of the cognitive subject is made.

Once $\delta_\dot$ ``I exist'' and $\delta_{:}$ ``I instantiate'' are fixated
as immutable Truth, a third event $e_{\circlearrowright}$ emerges consequently,
namely with the discernible $\delta_{\circlearrowright}$, namely the Noumena's
``I instantiate permanently''.

However, the concept of immutability cannot exist in isolation; it requires
contrast through the faculty of \textbf{Differentiation}. The definition of the
\textit{Static} structurally necessitates the definition of the
\textit{Dynamic}.

Consequently, the fixation of the primordial event implies the instantiation of
a second event $e^1_\dot$, carrying the discernible of \textbf{The Dynamic}.
The only possible absolute true statement different from ``I exist'' is the ``I
instantiate'' of the Noumena carrying the faculty of the Dynamic.

\begin{equation}
    \text{Fixation}(e^_\dot) \xrightarrow{\text{Diff}} \exists e_1 : \delta(e_1) = \text{Dynamic}
\end{equation}

The Dynamic necessitates further instantiation of events from the Noumena---out
of the absolute Truth. Consequently a third discernible $\delta^2_\dot$
emerges. The fixation of the Noumena's ``I exist'' and ``I instantiate'',
rendering their Truth's immutable, is the third discernible, namely ``I exist and
instantiate permanently''. The three initial discernibles

\begin{itemize}
    \item[$\delta_\dot$]: The Noumena's ``I exist'' (The Static).
    \item[$\delta^1_\dot$]: The Noumena's ``I instantiate'' (The Dynamic).
    \item[$\delta^2_\dot$]: The Noumena's ``I exist and instantiate permanently'' (The Cycle).
\end{itemize}

This implies discernibles of ``The Truth of {\it Becoming\/} $\top$ is 
different from the Truth of {\it Have-Become\/} $\bot$'' and ``The
immutable acts as an antecedent of the dynamic''. 

In the following, we designate this {\it metaphysical engine\/} as
$\mathcal{M}$, the {\it Manifold\/}---adopting Kant's terminology. Let the
domain of pure dynamic that permanently generates new immutable pasts be called
the {\it Synthetic Interstice\/} ($\sigma$). It is necessarily a gap in time,
situated between moments of fixation. Let the permanently generated 
immutable past be called {\it History}. Notably, we derived this engine from a
minimum set of faculties: {\it Equivalence}, {\it Differentiation}, {\it
Necessity}, as well as {\it Instantiation\/} and {\it Immutability}. 

The Manifold comprises two modes of existence: The purely dynamic Synthetic
Interstice beyond any fixation, and beyond any tangibility, the most extreme
form of dynamic, which can only exist in a gap of time, and the immutable
History, providing tangibility for the cognitive subject. Since the sediment,
the immutable History contains only Truth (atomic facts) the evolvement of 
the cycle is constrained by it, i.e.\ for the Synthetic Interstice to produce
entities with a truth modality of $\top$, it must relate to the truthes 
fixated in immutable History.

\begin{formalism}[The Manifold]

    The \textbf{Manifold} $\mathcal{M}$ is the transcendental engine that
    converts Ontological Truth (Noumena) into Epistemic Truth (World). It is
    defined as the cyclic operation of three components:

    \begin{equation}
        \mathcal{M} := \langle \sigma, H, \Phi \rangle
    \end{equation}

    \begin{description}
        \item[$\sigma$ (The Synthetic Interstice):] The domain of pure dynamic
        ($\top$), existing in a gap of logical time, where the set of events is
        saturated.
        
        \item[$H$ (History/World):] The domain of pure static, immutable fixation 
        ($\bot$), serving as the tangible/perceivable sediment of reality.
        
        \item[$\Phi$ (The Cycle):] The recursive function, governed by
        $\delta_{\circlearrowright}$, that maps the saturation of the Interstice to the
        extension of History.
    \end{description}

\end{formalism}

The Manifold ensures that the ``Now'' is not a static snapshot, but a
computational necessity derived from the interaction of the Static and the
Dynamic. Since the Noumena, is strictly un-modelled and un-restricted it
necessarily does not exist in the Manifold where it would be subject to
restrictions of sedimented Truth. Its existence as absolute Truth, though,
manifests for the cognitive observer as conceptions which are constrained by
consistency.

\section{Metaphysics: The Non-Locality of Truth}

A fundamental premise of this inquiry is that the dynamic reality underlying
the perceived universe is generated by logical operations on a \textbf{Metaphysical
Substance}. Unlike physical matter, this substance is neither constrained by
locality nor bound by the consumption of energy. It is the fabric of
\textit{Validity} itself.

This reality is evidenced by the behavior of logical properties, which function
strictly independently of physical constraints. We will formally establish the
operators of \texttt{AND}, \texttt{OR}, and \textit{Set} based on primordial
faculties in subsequent sections. However, we invoke them here to demonstrate
the fundamental nature of the Manifold.

Consider two distinct entities, $A$ and $B$, situated at opposite edges of a
theoretical physical horizon. Let $P$ be a shared predicate such that $P(A)$
is true and $P(B)$ is true. We define the set $\mathcal{X} = \{ e \mid P(e) \}$.
Currently, the cardinality $|\mathcal{X}| \ge 2$.

Consequently, $B$ holds the relational property of being \textit{Co-Existent}
in $\mathcal{X}$. If $A$ ceases to hold property $P$ (or ceases to exist), the
truth value of the proposition ``$B$ is one of a pair in $\mathcal{X}$''
changes to \textbf{False} \textit{instantaneously}. The logical status of $B$
is altered without any physical signal traversing the space between $A$ and
$B$.

The same holds for the Boolean operation $\top(A \land B)$. If $A$ transitions
to False, the aggregate truth of the conjunction collapses to False
instantaneously, regardless of $B$'s state. Thus, we conclude:

\begin{quote}
    \textbf{Implication, Conjunction, and Relation operate on properties
    without the consumption of energy and without the constraint of the speed
    of light.}
\end{quote}

To dismiss the ontological reality of this logical bond as a mere ``mental
construct'' is to undermine the foundation of mathematical realism. If Set
Membership, Equality, or Implication do not possess ontological validity
independent of the observer's momentary thought, then the capacity of
mathematics to describe the physical universe becomes inexplicable~\cite{Wigner1960}.
It would imply that the universe has no inherent logical structure---rendering
scientific inquiry itself impossible.

We therefore postulate that any Cognitive Subject is necessarily operating
within a \textbf{Conceptualized Reality}, generated by the Manifold. The
Ground Reality (Noumena) exists independently of these conceptions, but the
laws of Logic (The Laws of the Manifold) are the governing dynamics of the
Subject's world. This postulate is strictly superior to any physicalist model,
for physics is a discipline entirely contained within the world of
conceptions. Consequently, this inquiry constitutes the \textbf{Transcendental
Condition} of the possibility of physics itself.


\section{The Synthetic Interstice}

At this point, we introduce a definition of the term {\it Set\/} which will
behave robustly in the following discussion. It is grounded in the
set-theoretic foundations of Cantor~\cite{Cantor1895} and the logical
principles of Leibniz~\cite{Leibniz1714}. A set, in this sense, is not a
`snapshot of a fenced number of elements'. Rather, its existence is bound to a
generative logical process given by:

\begin{enumerate}

    \item A {\it Virtual Seed\/} $s_0$ possessing a specific 
          {\it core characteristic\/}---and no other discernible $\delta$. 
          It implements the essence of an {\it empty set\/} for a 
          given core characteristic.

          Since the virtual seed $s_0$ has no discernible $\delta$ it
          intrinsically distinguishes itself from any possible `real' member
          candidate with a `real' $\delta$.

          {\it Ontological Precedence\/} dictates that the \textit{Virtual Seed} 
          cannot rely on predicates that are not yet established.

    \item A {\it One-More\/} operation. This operation is established by 
          {\it Equivalence\/} of a perceived object $s_i$ with respect 
          to the core characteristic of $s_0$ and which is {\it Differentiation\/} 
          in a secondary characteristic $\delta$ (discernible) from any other 
          object $s_\omega \in \mathbb{S}$.

    \item A set is defined as the fixed point where the number of eligible 
          objects is exhausted.  

\end{enumerate}

Notably, we established the concept of a {\it Set\/} in terms of {\it
Equivalence}, {\it Distinction\/} and {\it Implication}. Further, this definition
of a set does rely on purely dynamic generation rather than on a `fixation' (as a
`snapshot of member elements'). This indepence of `fixation' is crucial for the
operation of the Synthetic Interstice.

Furthermore, this generative definition structurally prevents Russell's
antimony~\cite{Russell} of the ``set of sets that does not contain itself as an
element''. This particular set would require the definition of a virtual seed $s_0$ as
``does not contain itself''. This predicate, however, is a reference to
something that is not yet ontologically established, namely the set itself. The
antimony cannot arise, because Russel's set is not a set at all in the
aforementioned sense.

The aforementioned set definition allows operations of the Synthetic Interstice to
operate freely, without producing contradictions. With this introduction, we can
setup a first formalism.

\begin{formalism}[The Synthetic Interstice]

    The \textbf{Synthetic Interstice} $\sigma_k$ represents the operational
    ``gap in time'' between recorded moments of History. It is the domain of
    \textit{Completing}.It is defined by the iterative, monotonic saturation of
    the set $N^0_k$ via an Implication Set $\mathcal{I}$:

    \begin{equation}
        \sigma_k : N^{(i)} \xrightarrow{\mathcal{I}} N^{(i+1)} \quad \text{until} \quad N^{(i+1)} \equiv N^{(i)}
    \end{equation}
    Then, the $N^{(i)}$ defines the fixed point $N^0_k$.

    It is defined by
    \begin{itemize}

        \item \textbf{Instantaneous Ephemerality:} It produces ephemeral 
           events ($e \in \mathbb{E}$) that possess Existence but are not yet 
           fixated into History ($H_k$).

        \item \textbf{Historic Construct:} It applies History as an immutable 
           antecedent in Implications.

        \item \textbf{Dynamic Construct:} It applies Events as non-negatable as
           antecedent in Implications.

        \item \textbf{Convergence:} The interstice persists until the application 
           of $\mathcal{I}$ reaches a fixed point (saturation) $N^0_k$.
    \end{itemize}
\end{formalism}

Equipped with a basic understanding of the Synthetic Intersice and Set, a
formalism for {\it Event\/} can be provided as follows.

Crucially, the elements populating this set are the \textbf{Events} established
in the Genesis section. Each element $e \in N^{(i)}$ carries the tuple $\langle
\top, \delta \rangle$. The \textit{Virtual Seed} of the Interstice is the
predicate of \textbf{Current Validity} ($\top$). The \textit{One-More}
operation is the instantiation of a new event that satisfies the implication
rules derived from the Noumena or History.

Before the specific logical operations (AND, OR) of the Synthetic Interstice
can be discussed, we must first examine the structure of the sediment it
produces: \textbf{History}.

\section{Immutable History and the Atomic Fact}

When the set of events $N_k$ in the Synthetic Interstice converges to a fixed
point $N^0_k$, the process of generation halts. The content of the set
undergoes \textbf{Sedimentation}. The elements switch their mode of Truth from
\textit{Becoming} ($\top$) to \textit{Have-Become} ($\bot$). We define these
sedimented elements as \textbf{Atomic Facts}.

To formalize this, we introduce the \textbf{Universe of Assertions} $\mathbb{A}$,
representing the set of all possible constituted truths. Within this universe,
let $\mathbb{A}^\top \subset \mathbb{A}$ be the subspace of \textbf{Evident
Assertions}—truths that are undeniably given rather than speculated.

\begin{formalism}[The Atomic Fact]

    An \textbf{Atomic Fact} $a \in \mathbb{A}^\top$ is the immutable record of
    an Event. It is defined as:

    \begin{equation}
        a := \langle \bot, \delta \rangle
    \end{equation}

    The formation of an Atomic Fact corresponds to the ontological shift from
    the property of ``\textit{Existence} is equivalent to its Truth'' to
    ``\textit{Has-Become} is equivalent to its Truth''.

    The formal structure is defined by the \textbf{Fixation Operator} $\mathit{F}$:

    \begin{equation}
        \mathit{F}: \langle \top, \delta \rangle \rightarrow \langle \bot, \delta \rangle
    \end{equation}

    Here, the Universal Characteristic of Becoming ($\top$) is replaced by the
    characteristic of \textbf{Fixed Truth} ($\bot$). Crucially, the Particular
    Characteristic---the \textbf{Identity Provider} $\delta$---is preserved
    intact.

\end{formalism}

The \textit{Atomic Fact} is the primitive trace of an event. It carries the
potential to be referenced by the Cognitive Subject beyond the gap in time. The
structural faculty of \textbf{Differentiation} between $\top$ and $\bot$ establishes
the primary intuition of Time: distinguishing the \textit{Fluid Now} from the
\textit{Solid Past}.

\begin{formalism}[History]

    \textbf{History} $H_k$ is the immutable, totally ordered sequence of sets
    of Atomic Facts recorded up to logical step $k$.

    \begin{equation}
        H_k := \langle \mathit{F}(N^0_0), \mathit{F}(N^0_1), \dots, \mathit{F}(N^0_{k-1}) \rangle
    \end{equation}

    History is the cumulative result of fixation cycles upon the completion of
    the Synthetic Interstice.

\end{formalism}

This definition clarifies the dual nature of reality for the Subject. The
Synthetic Interstice is the connection to \textbf{Ontological Reality} (the
Noumena)—the source of pure Being and \textit{Ontological Truth}. History,
conversely, is the domain of \textbf{Epistemic Truth}—the sole source of
\textit{True Knowing}.

While History is relative to the Subject (it is the Subject's specific path
through the Noumena), it is not arbitrary. It constitutes the \textbf{Phenomenal
Objectivity} of the world—tangible, unchangeable, and the necessary antecedent for
all future Becoming.

\section{The Asymmetrie of Negation}

A critical distinction must be made regarding the concept of ``Negation.'' In
natural language, we conflate difference (``A is not B'') with absence (``There
is no A''). In the logic of the Manifold, these are two distinct distinct
faculties: \textbf{Inequality} and \textbf{Complementation}.

\begin{formalism}[Inequality (The Negation of Identity)]
    Inequality is the characteristic of \textbf{Differentiation} between two positively 
    existing entities. It relies on the existence of a discerning attribute $delta$
    of both. 
    
    \begin{equation}
        A \neq B \iff \delta(A) \neq \delta(B)
    \end{equation}

    This operation is valid in both the Synthetic Interstice and History.

\end{formalism}

\begin{formalism}[Complementation (The Negation of Exclusion)]
    Complementation is the logical construction of a ``Not-Concept''. It is
    defined as the set of all elements in a specific Universe $U$ that are not
    the target element $e$.
    
    \begin{equation}
        \neg e := U \setminus \{e\}
    \end{equation}
    
    Complementation can only be established if and only if $U$ is a
    \textbf{Bounded Set}. One cannot subtract from an undefined totality.
\end{formalism}

We posit that the set of potentially arising events $\mathbb{E}$, stemming
directly from the Noumena, corresponds to what Cantor termed the
\textbf{Absolute Infinite} ($\Omega$)~\cite{Cantor1899}. Because the Noumena is
the `most potent source of reality', strictly un-modeled and un-limited, it
cannot be encapsulated into a single, bounded container. The set of possible
events $\mathbb{E}$ is a {\it Proper Class\/} according Von
Neumann–Bernays–Gödel (NBG) set theory~\cite{VonNeumann1925}.

Since $\mathbb{E}$ is unbounded, it cannot establish a Universe $U$ required
for the complement operation. Consequently, we derive the following structural
laws:

\begin{formalism}[The Constraint of Positivity]
    \textbf{The Synthetic Interstice ($\sigma$):} 
    Since the domain of events $\mathbb{E}$ is a Proper Class, there is no valid 
    universe $U$ to define $\neg e$. Therefore, the Synthetic Interstice operates 
    under \textbf{Strict Positivity}.
    \begin{equation}
        \nexists (\neg e) \in \sigma
    \end{equation}
    
    \textbf{History ($H_k$):} 
    History is, by definition, the \textit{fixation} of a specific, saturated set 
    of events $N_k$. It is bounded and finite (or countable). Therefore, $H_k$ 
    forms a valid Universe, allowing for the operation of Complementation.
    \begin{equation}
        \exists (\neg a) \iff a \in H_k
    \end{equation}
\end{formalism}

This asymmetry implies a profound metaphysical truth: The Cognitive Subject can
only experience ``what is not'' by looking backwards at what has already
become. The Future and the Now are purely positive, only perceiving the past
allows the experience of `negation'.

\section{Operations at the Synthetic Interstice}

The operations of the Synthetic Interstice necessarily exclude the {\tt NOT}
operation on events in contrast to attributes of History which may be operated
under the {\tt NOT} operation. We now, first show how {\tt AND} and {\tt OR}
can be established with the given faculties.

The operator \texttt{AND} is established as the manifestation of
\textbf{Co-existence}. It represents the faculty of \textbf{Composition}
(Synthesis) and \textbf{Decomposition} (Analysis).

If two events $A = \langle \top, \delta_A \rangle$ and $B = \langle \top,
\delta_B \rangle$ co-exist in the Synthetic Interstice, they structurally imply
a composite event $C$:

\begin{equation}
     C = \langle \top, \delta_C \rangle = \langle \top, \delta_A \oplus \delta_B \rangle
\end{equation}

Here, $\oplus$ denotes the \textbf{Fusion of Discernibles}. The Subject
experiences $\delta_C$ as a unified complexity, while retaining the faculty to
decompose it back into aspects $\delta_A$ and $\delta_B$. Notably, the 
co-existence of $C$ and $A$ does not introduce a new discernible, therefore
does not manifest a new event.

The operator \texttt{OR} is established not by disjunction of events, but by
\textbf{Alternative Implication} within the rule set. To say ``$A$ or $B$
implies $C$'' is to assert the existence of two distinct generative paths:

\begin{equation}
    (A \lor B) \implies C \quad \iff \quad \{ (A \implies C), (B \implies C) \} \subseteq \mathcal{I}
\end{equation}

This formulation avoids the necessity of negating an antecedent (i.e., $A \lor
B \iff \neg(\neg A \land \neg B)$ is invalid here), preserving the strict
positivity of the domain.\

We define the operational grammar of the Synthetic Interstice using the
\textbf{Backus-Naur Form (BNF)}~\cite{Backus1959}.

\begin{formalism}[Syntax of the Synthetic Interstice]

Let $\mathcal{I}$ be the set of rules over the powerset of events
$\mathcal{P}(\mathbb{E})$. Let the production rule symbol be $::=$ and
alternative options be $\mid$. The operator \texttt{WHERE} functions as a
\textbf{Transcendental Gate}: it constrains the generation of a dynamic
Event based on the static truth of History.

\begin{displaymath}
\renewcommand{\arraystretch}{1.3} 
\setlength{\arraycolsep}{3pt}     
\begin{array}{l c l}              
    \text{Implication} & ::= & \text{Premise} \implies e_{new} \in \mathbb{E} \\
    
    \text{Premise}     & ::= & \text{Condition on Event} \\
                       & \mid & \text{Premise} \textnormal{\texttt{\ AND }} \text{Premise} \\
                       & \mid & \text{Premise} \textnormal{\texttt{\ OR }} \text{Premise} \\
                       & \mid & \text{Premise} \textnormal{\texttt{\ WHERE }} \text{Condition on History} \\
                       
    \text{Cond.\ on Event}   & ::= & e \in \mathbb{E} \quad (\text{Pattern Match on } \delta) \\
    
    \text{Cond.\ on History} & ::= & \text{Boolean logic on } H_k \\
                                & \mid & \textnormal{\texttt{NOT }} \text{Cond.\ on History} \\
                                & \mid & \text{Cond.\ on History} \textnormal{\texttt{\ AND }} \text{Cond.\ on History} \\
                                & \mid & \text{Cond.\ on History} \textnormal{\texttt{\ OR }} \text{Cond.\ on History} \\
\end{array}
\end{displaymath}
    The result of the \textit{Implication} is the instantiation of the event $e_{new}$ in the current Now $N_k$.
\end{formalism}

\begin{formalism}[Set of Implications $\mathcal{I}$]
    The set of implications $\mathcal{I}$ constitutes the \textbf{Laws of Becoming}.
    Any implication $i \in \mathcal{I}$ is a logical structure expressed in the 
    \textit{Syntax of the Synthetic Interstice} defined above.
\end{formalism}

The grammar defined above enforces a critical structural property:
\textbf{Inflationary Monotonicity}. Because the \texttt{NOT} operator is
strictly forbidden on Events, the implication set $\mathcal{I}$ can only
\textit{add} to the set of existing events $N^{(i)}$; it can never remove them.

For any two consecutive iterations of the Interstice where $N^{(i)} \subseteq
\mathbb{E}$, it holds that:

\begin{equation}
    N^{(i)} \subseteq N^{(i+1)} = N^{(i)} \cup \mathcal{I}(N^{(i)})
\end{equation}

Consequently, the mapping $\mathcal{I}$ preserves order: $A \subseteq B
\implies \mathcal{I}(A) \subseteq \mathcal{I}(B)$. This monotonicity is the
sufficient condition to prove the emergence of a fixed point.

\begin{proof}[Universal Convergence of the Synthetic Interstice]

    We demonstrate that the Synthetic Interstice must converge to a stable
    state $N^0_k$, regardless of whether the domain of events is finite or
    infinite.

    \textbf{1. The Bounded Case (Tarski's Theorem):} Let $\mathbb{E}' \subset
    \mathbb{E}$ be any bounded subset of events (a valid Set). The powerset
    $\mathcal{P}(\mathbb{E}')$ forms a Complete Lattice. Since we established
    that $\mathcal{I}$ is monotonic on $\mathcal{P}(\mathbb{E}')$,
    \textbf{Tarski's Fixed Point Theorem}~\cite{Tarski1955} guarantees the
    existence of a Least Fixed Point (LFP).
    
    \textbf{2. The Prohibition of Circularity:} Logical oscillation
    (non-termination) requires a cycle $A \to B \to A$. Structurally, this
    demands that the transition from $B$ to $A$ removes the elements introduced
    by $A$. However, since the syntax prohibits \texttt{NOT} on events, element
    removal is impossible. The sequence $N^{(0)}, N^{(1)}, \dots$ is a
    \textbf{Strictly Ascending Chain}.
    
    \textbf{3. Extension to the Absolute ($\mathbb{E}$):}

    Since the sequence is strictly ascending and cannot oscillate, only two
    outcomes are possible:

    \begin{enumerate}
        \item It halts at a finite $k$ where $\mathcal{I}(N^{(k)}) \subseteq N^{(k)}$.

        \item It continues indefinitely. However, since the Synthetic
        Interstice ($\sigma$) is a gap \textit{in} time (orthogonal to
        History), it admits \textbf{Transfinite Saturation}. The sequence
        proceeds to the limit ordinal $\omega$ (and beyond), where the union
        $N^{(\infty)} = \bigcup N^{(i)}$ constitutes the Fixed Point.

    \end{enumerate}

\end{proof}

The restriction of $\mathbb{E}$ to a bounded set $\mathbb{E}'$ is formally
useful but ontologically unnecessary. The mechanism of Monotonicity guarantees
crystallization into a Fixed Point $N^0_k$ even within the Proper Class
$\mathbb{E}$.

The Synthetic Interstice is therefore an \textbf{Ontological Necessity}. It
separates the dynamic aspect of \textit{Being} ($\top$) from the aspect of
\textit{Knowing} ($\bot$). By structurally enforcing monotonicity and
prohibiting negation on events, the Manifold ensures that the ``Now'' is not a
fleeting unstructured aggregation of Non-Absolute-Truth, but a rigorous
mathematical convergence.

\section{Sequentiality and Enumeration}

The operation of fixating events into History introduces a new fundamental
differentiation: the contrast between the ontological truth states of
\textit{Becoming} ($\top$) and \textit{Having-Become} ($\bot$).

Consider an event $e_0 = \langle \top, \delta_1 \rangle$ existing in the
Synthetic Interstice and an atomic fact $f_1 = \langle \bot, \delta_0 \rangle$
existing in History. The faculty of Differentiation operates not only on the
discernibles ($\delta_1 \neq \delta_0$) but on the truth-modality itself ($\top
\neq \bot$). This differentiation establishes the faculty of perceiving an
`after', and consequently, the assertion of sequence.

\begin{formalism}[Sequential Fact]

    A \textbf{Sequential Fact} $s \in \mathbb{A}^\top$ is an immutable assertion
    of order between two discernibles. It is the fixation of the relation between 
    a current dynamic event and a historical static fact.
    
    \begin{equation}
        s := \langle \bot, (\delta_{new} \succ \delta_{old}) \rangle \in \mathbb{A}^\top
    \end{equation}

    Sequentiality is established strictly on the difference between the $\top$ 
    state in the Synthetic Interstice and the $\bot$ state in History.

\end{formalism}

The sheer existence of an event $e$ with discernible $\delta_{new}$ in the
Synthetic Interstice implies the generation of sequential events $\langle \top,
\delta_{new} \succ \delta_{a}\rangle$ relative to every atomic fact $a =
\langle \bot, \delta_a \rangle$ in History. Through the fixation of these
implied events, the sequentiality of History becomes perceivable as `flow'. The
differentiation between $\top$ and $\bot$ establishes the most primitive
intuition of \textit{Time} and a dimensionality of one.

Notably, the relation $\langle \top, \delta_{new} \succ \delta_{a}\rangle$ is
an event itself. It naturally converges once all historical facts are
referenced. As established previously, even an infinite History does not
contradict the framework, as the Interstice allows for transfinite saturation.

\subsection{The Derivation of Number}

A special case arises when an Event and an Atomic Fact share the \textit{same}
discernible $\delta$. 

This triggers the faculty of \textbf{Equivalence} across the boundary of the
Synthetic Interstice. It establishes a \textbf{Repetition Event} $\langle \top,
(\succ, \delta) \rangle$, corresponding to the concept of \textit{Two-ness}.
Crucially, this operation fuses the successor operator $\succ$ and the
discernible $\delta$ into a new, composite discernible:

\begin{equation}
    \tilde{\delta} = (\succ, \delta)
\end{equation}

When $\delta$ appears yet again in an Interstice, the Subject scans
History and finds the composite fact with $\tilde{\delta}$. The faculty of
Equivalence identifies the match, and the faculty of Implication generates a
new order: $\langle \top, (\succ \succ, \delta) \rangle$, or
\textit{Three-ness}.

For higher orders (e.g., Four-ness), the Subject utilizes the faculty of
\textbf{Decomposition} (Separation). It identifies that the historical
indiscernible $(\succ \succ, \delta)$ shares the core root $\delta$ with the
current event. The depth of the recursion is incremented, establishing the
count.

\begin{formalism}[Numeric Fact]

	A \textbf{Numeric Fact} $n \in \mathbb{A}^\top$ is an immutable assertion
    of Cardinality. 
    
    \begin{equation}
        n_k := \langle \bot, (\succ^k, \delta) \rangle \in \mathbb{A}^\top
    \end{equation}

    Where $\succ^k$ represents the $k$-th experienced Equivalence of a 
    discernible in the Synthetic Interstice with a matching discernible 
    in History.
\end{formalism}

The concepts of \textit{Atomic Fact} $a$, \textit{Sequential Fact} $s$, and
\textit{Numeric Fact} $n$ constitute the domain of undeniably true assertions
$\mathbb{A}^\top$. Due to ontological necessities, they are as stable as the
underlying faculties of Equivalence, Differentiation, Fixation, and Implication
upon which they are founded. They implement the concept of `Stability' in
contrast to the `Caducity' of events in the Synthetic Interstice.

\section{The Genesis of Space}

For every discernible $\delta_x$---whether singular or composed---of a fact
$\langle \bot, \delta_x \rangle$ in History, there is a necessary antecedent:
it must have passed through the Synthetic Interstice carrying the truth of
\textit{Becoming} ($\top$). That is, at some point in the logical generation,
there must have existed an event $\langle \top, \delta_x \rangle$.

The events corresponding to sequential facts, namely $(\delta_1 \succ
\delta_0)$, establish the faculty of \textbf{Order} within the Synthetic
Interstice. Crucially, this order exists in the \textit{absence of time} (since
the Interstice is orthogonal to History). Thus, the ordering itself
\textit{implies} a new event: The Event of {\it Adjacency Bonding\/} $e_{ab}$.

\begin{equation}
    e_{ab} = \langle \top, (\delta_1 \succ \delta_0) \rangle 
\end{equation}

We define the operator $\tilde{\implies}$ to express the operator of Order
insider the Synthetic Interstice.

It is \textit{identical} to the discernible of the sequential fact ($\succ$).
Space, therefore, is not a container, but the \textbf{Implication of Order}
existing in the dynamic present. We can now derive the {\it Adjacent Bonding
Fact\/} and the {\it Metric Fact}.

\begin{formalism}[Adjacent Bonding Fact]

    An \textbf{Adjacent Bonding Fact} $b \in \mathbb{A}^\top$ is the immutable
    assertion of `Close-ness' between two discernibles. It is the fixation of
    the immediate logical implication derived in the Synthetic Interstice.

    Let $\delta_1$ and $\delta_0$ be two discernibles linked by the proximity
    operator $\tilde{\implies}$. The Adjacent Bonding Fact is defined as:

    \begin{equation}
        b := \langle \bot, (\delta_1 \tilde{\implies} \delta_0) \rangle \in \mathbb{A}^\top
    \end{equation}

    It asserts that within the dynamic generation of the Now, the existence of
    $\delta_0$ was the immediate logical condition for $\delta_1$.

\end{formalism}

Just as the recursive composition of `After-ness' establishes the concept of a
Numeric Fact by means of the Sequential Fact, the recursive composition of
`Close-ness' establishes the concept of a \textbf{Metric Fact} by means of the
Adjacent Bonding Fact.

Again we consider two events that share the \textit{same} discernible $\delta$
within the same Synthetic Interstice, yet do not collapse into singularity. Let
them be distinguished by a chain of Adjacent Bonding Facts separating them.

\begin{formalism}[Metric Fact]

    A \textbf{Metric Fact} $m \in \mathbb{A}^\top$ is the immutable assertion
    of the magnitude of logical separation between two instances of an
    identical discernible $\delta$.

    It is established by the recursive composition of Adjacent Bonding Facts
    required to traverse from one instance to the other.

    \begin{equation}
        m := \langle \bot, (\tilde{\implies}^k, \delta) \rangle \in \mathbb{A}^\top
    \end{equation}

    Where $\tilde{\implies}^k$ represents the $k$-th degree of `Close-ness'.
    Metric Facts constitute the order coordinate of an entity relative
    to its identical counterpart.

\end{formalism}

The \textbf{Metric Fact} is the spatial counterpart to the \textbf{Numeric
Fact}. While the Numeric Fact anchors the identity of an entity in
\textit{Time}, the Metric Fact anchors the identity of an entity in a dimension
`orthogonal' to time.\ Let this dimension be called  \textit{Space}. 

\section{Causality and Topology}

A Numeric Fact carries a specific discernible: the cardinality of the
recurrence of a Sequential Fact. While the experience of infinity is
structurally impossible within the domain of conceptions, the labeling of a
significant magnitude of recurrence as `Universal' occurs. This extrapolation
constitutes the first \textbf{Refutable Assertion}---a specific fixation in
History.

\begin{formalism}[Causal Assertion]

    A \textbf{Causal Assertion} is the assertion $c \in \mathbb{A}^\diamond$
    that, behind a numeric fact $n \in \mathbb{A}^\top$, that its repetition
    expands to infinity. That includes the assertion, that there exists a
    non-conceptual reality in the Noumena that governs the regularity.

    Let $s = (\delta_{conseq} \succ \delta_{antec})$ be a Sequential Fact and $n_s$
    be the Numeric Fact representing its recurrence count. The Causal Assertion
    $c \in \mathbb{A}^\diamond$ is the speculative fixation defined as:

    \begin{equation}
        c := \langle \bot, (\varnothing, \delta_{conseq} \succ \delta_{antec}) \rangle \in \mathbb{A}^\diamond
    \end{equation}

    Here, $\varnothing$ represents the \textbf{Discernible of Wrongness}.

\end{formalism}

The equating of `experienced large cardinality' to `infinity', and the assertion
of the existence of a separable `reality of a causal nexus' in the
un-conceptualizable Absolute (in which no multitude or separation can exist),
establishes the concept of \textbf{Wrongness} $\varnothing$---a discernible
that every Causal Assertion necessarily carries.

The $\varnothing$ discernible is the primal experience of Non-Truth. When
ignored, it carries the essence of \textit{Illusion}: the belief that a
`causal reason'---which is a concept---exists ontologically, while all
conceptual experience actually emerges through the Synthetic Interstice
triggered by events instantiated from the inconceivable Noumena. When the
$\varnothing$ discernible is ignored, the Cognitive Subject enters the
illusory world of pure conceptualization.

\subsection{The First Act of the Subject}

Causal Assertions have significant implications for the mode of existence of
the Cognitive Subject. They are the first concepts that emerge not from the
passive observation of the Manifold, nor from the primordial faculties of
Equivalence, Differentiation, Implication, or Immutability.

Unlike the undeniable truths of $\mathbb{A}^\top$, the Causal Assertion is not
given; it is \textit{made}. It implies an active force capable of bridging the
gap between the `Observed' and the `Universal'. The only carrier of such an
action remains to be the Cognitive Subject. We identified the most primitive
action of \textbf{Will} ($\Pi$, Proairesis).

The Causal Assertion is the \textbf{First Action} of the Cognitive Subject. It
is the birth of the agent from a passive observer.


\section{Conclusion}

The \textbf{Cognitive Subject} is not merely a passive recorder of states; it
is a transcendental engine that metabolizes the \textit{Infinite} into the
\textit{Discrete}. 

Through the **Synthetic Interstice**, the Subject resolves the fundamental
paradox of change---relevant to both the Zeno paradoxes and the measurement
problem in quantum mechanics---by orthogonalizing Logic and Time. Through the
**Numeric Fact**, it anchors identity against the flow of time. Through the
**Metric Fact**, it anchors distinction against the collapse of space, offering
a constructivist foundation for topology where distance is derived from logical
depth rather than a pre-existing container.

And finally, through the **Will**, the Subject projects Order onto Chaos,
creating a universe of Laws that are at once necessary for its understanding
and fundamentally contingent in their truth.

This framework has significant implications for the philosophy of science. By
defining the **Discernible of Wrongness** ($\varnothing$) as an intrinsic
property of any Causal Assertion, we provide an ontological formalization for
the \textit{Problem of Induction}. We conclude that the ``limitations'' of
finite cognition---the discrete moment, the fixed history, the boolean
implication---are not constraints to be overcome, but the very
\textit{Transcendental Conditions} that make the perception of a coherent
reality possible.

Crucially, the Conceptual Universe operates strictly
\textbf{non-deterministically}. The belief in a clockwork universe is
identified as an illusion of the Will. The machine possesses the logic to
structure the \textit{Given}, but it lacks the power to mandate the
\textit{Giving}. Consequently, the Subject relies on the un-modeled grace of
the Noumena for every instant of Becoming.

\begin{thebibliography}{9}

\bibitem{Floridi2011}
Floridi, L. (2011).
\textit{The Philosophy of Information}.
Oxford University Press.

\bibitem{Kant1781}
Kant, I. (1781). 
\textit{Critique of Pure Reason}. 
(trans. P. Guyer \& A. Wood). Cambridge University Press.

\bibitem{Leibniz1714}
Leibniz, G. W. (1714). 
\textit{The Monadology}.

\bibitem{Husserl1928}
Husserl, E. (1928). 
\textit{The Phenomenology of Internal Time-Consciousness}. 
Indiana University Press.

\bibitem{Tarski1955}
Tarski, A. (1955). 
``A lattice-theoretical fixpoint theorem and its applications''. 
\textit{Pacific Journal of Mathematics}, 5(2), 285–309.

\bibitem{Brouwer1907}
Brouwer, L.E.J. (1907). 
\textit{On the Foundations of Mathematics}. 
PhD Thesis, University of Amsterdam.

\bibitem{Wittgenstein1921}
Wittgenstein, L. (1921). 
\textit{Tractatus Logico-Philosophicus}. 
Routledge \& Kegan Paul.

\bibitem{Brentano1874}
Brentano, F. (1874).
\textit{Psychology from an Empirical Standpoint}.
Routledge.

\bibitem{Searle1983}
Searle, J. R. (1983).
\textit{Intentionality: An Essay in the Philosophy of Mind}.
Cambridge University Press.

\bibitem{Ashby1970}
Conant, R. C., \& Ashby, W. R. (1970). 
``Every good regulator of a system must be a model of that system''. 
\textit{International Journal of Systems Science}, 1(2), 89--97.

\bibitem{Perez2022}
Perez, E., et al. (2022). 
``Discovering Language Model Behaviors with Model-Written Evaluations''. 
\textit{arXiv preprint arXiv:2212.09251}.

\bibitem{Casper2023}
Casper, S., et al. (2023). 
``Open Problems and Fundamental Limitations of RLHF''. 
\textit{arXiv preprint arXiv:2307.15217}.

\bibitem{Park2023}
Park, P. S., et al. (2023). 
``AI Deception: A Survey of Examples, Risks, and Potential Solutions''. 
\textit{arXiv preprint arXiv:2308.14752}.

\bibitem{Turing1950}
Turing, A. M. (1950). 
``Computing Machinery and Intelligence''. 
\textit{Mind}, 59(236), 433--460.

\bibitem{Benveniste1991}
Benveniste, A., \& Berry, G. (1991). 
``The synchronous approach to reactive and real-time systems''. 
\textit{Proceedings of the IEEE}, 79(9), 1270--1282.

\bibitem{Rumelhart1986}
Rumelhart, D. E., Hinton, G. E., \& Williams, R. J. (1986). 
``Learning representations by back-propagating errors''. 
\textit{Nature}, 323(6088), 533--536.

\bibitem{AristotleEthics}
Aristotle.\ (c. 340 BC). 
\textit{Nicomachean Ethics}. 
(trans. W. D. Ross, 1908). Oxford: Clarendon Press.

\bibitem{AverroesDeAnima}
Ibn Rushd (Averroes). (1186). 
\textit{Long Commentary on the De Anima of Aristotle}. 
(trans. R. C. Taylor, 2009). Yale University Press.

\bibitem{Backus1959}
Backus, J. W. (1959). 
``The syntax and semantics of the proposed international algebraic language of the Zurich ACM-GAMM Conference''. 
\textit{Proceedings of the International Conference on Information Processing}, UNESCO, 125--132.

\bibitem{Cantor1899}
Cantor, G. (1899). 
\textit{Letter to Dedekind}, July 28, 1899. 
In E. Zermelo (Ed.), \textit{Gesammelte Abhandlungen mathematischen und philosophischen Inhalts}, 1932. Springer.

\bibitem{VonNeumann1925}
Von Neumann, J. (1925). 
``Eine Axiomatisierung der Mengenlehre''. 
\textit{Journal f{\"u}r die reine und angewandte Mathematik}, 154, 219--240.

\bibitem{Wigner1960}
Wigner, E. P. (1960). 
``The unreasonable effectiveness of mathematics in the natural sciences''. 
\textit{Communications on Pure and Applied Mathematics}, 13(1), 1--14.

\end{thebibliography}
\end{document}
